\documentclass[a5paper,11pt]{ltjtbook}
\usepackage[pdfpagelayout=TwoPageRight,pdfdirection=R2L]{hyperref}
\begin{document}
\thispagestyle{empty}
\vspace*{3zw}
{\Huge\bfseries 本を読むための道具}\par\vspace{2zw}
{\Large 見本の文庫}\par\vspace{8zw}
{\large Book Viewer 見本}
\clearpage
\chapter{棚の奥の一冊}
夜の電車で数ページ、昼休みにまた数ページ。細切れの時間でも、前の続きからすぐに読み始められるなら、一冊を読み終えるのは思ったほど難しくない。本をスキャンしてデータにすると、この順序が崩れることがある。見開きの左右が入れ替わったり、白紙を一枚飛ばしたせいで、左右の組み合わせが一ページずれたりする。

文字の大きさや余白も、読みやすさを左右する。ページ全体を収めて眺めたいときもあれば、幅いっぱいに広げて一行ずつ追いたいときもある。雨の午後、棚の奥から古い文庫本を一冊取り出した。背は日に焼け、ページの端は少し波打っている。それでも開けば、活字は昔と同じ場所で待っていた。

紙の本を読むとき、私たちは右のページから左のページへ、上から下へと目を運ぶ。縦書きの本では、ページをめくる手も右から左へ動く。その身体の動きごと、読書という体験はできあがっている。紙の本を読むとき、私たちは右のページから左のページへ、上から下へと目を運ぶ。縦書きの本では、ページをめくる手も右から左へ動く。その身体の動きごと、読書という体験はできあがっている。

夜の電車で数ページ、昼休みにまた数ページ。細切れの時間でも、前の続きからすぐに読み始められるなら、一冊を読み終えるのは思ったほど難しくない。雨の午後、棚の奥から古い文庫本を一冊取り出した。背は日に焼け、ページの端は少し波打っている。それでも開けば、活字は昔と同じ場所で待っていた。

ずれた見開きは、読み進めるほどに気持ちが悪い。挿絵の右半分と左半分が別々のページに分かれてしまうこともある。一枚の空白ページを差し込むだけで、それはきちんと元に戻る。雨の午後、棚の奥から古い文庫本を一冊取り出した。背は日に焼け、ページの端は少し波打っている。それでも開けば、活字は昔と同じ場所で待っていた。

紙の本を読むとき、私たちは右のページから左のページへ、上から下へと目を運ぶ。縦書きの本では、ページをめくる手も右から左へ動く。その身体の動きごと、読書という体験はできあがっている。文字の大きさや余白も、読みやすさを左右する。ページ全体を収めて眺めたいときもあれば、幅いっぱいに広げて一行ずつ追いたいときもある。

文字の大きさや余白も、読みやすさを左右する。ページ全体を収めて眺めたいときもあれば、幅いっぱいに広げて一行ずつ追いたいときもある。紙の本を読むとき、私たちは右のページから左のページへ、上から下へと目を運ぶ。縦書きの本では、ページをめくる手も右から左へ動く。その身体の動きごと、読書という体験はできあがっている。

ずれた見開きは、読み進めるほどに気持ちが悪い。挿絵の右半分と左半分が別々のページに分かれてしまうこともある。一枚の空白ページを差し込むだけで、それはきちんと元に戻る。紙の本を読むとき、私たちは右のページから左のページへ、上から下へと目を運ぶ。縦書きの本では、ページをめくる手も右から左へ動く。その身体の動きごと、読書という体験はできあがっている。

文字の大きさや余白も、読みやすさを左右する。ページ全体を収めて眺めたいときもあれば、幅いっぱいに広げて一行ずつ追いたいときもある。雨の午後、棚の奥から古い文庫本を一冊取り出した。背は日に焼け、ページの端は少し波打っている。それでも開けば、活字は昔と同じ場所で待っていた。

\chapter{めくる方向}
紙の本を読むとき、私たちは右のページから左のページへ、上から下へと目を運ぶ。縦書きの本では、ページをめくる手も右から左へ動く。その身体の動きごと、読書という体験はできあがっている。ずれた見開きは、読み進めるほどに気持ちが悪い。挿絵の右半分と左半分が別々のページに分かれてしまうこともある。一枚の空白ページを差し込むだけで、それはきちんと元に戻る。

雨の午後、棚の奥から古い文庫本を一冊取り出した。背は日に焼け、ページの端は少し波打っている。それでも開けば、活字は昔と同じ場所で待っていた。文字の大きさや余白も、読みやすさを左右する。ページ全体を収めて眺めたいときもあれば、幅いっぱいに広げて一行ずつ追いたいときもある。

雨の午後、棚の奥から古い文庫本を一冊取り出した。背は日に焼け、ページの端は少し波打っている。それでも開けば、活字は昔と同じ場所で待っていた。ずれた見開きは、読み進めるほどに気持ちが悪い。挿絵の右半分と左半分が別々のページに分かれてしまうこともある。一枚の空白ページを差し込むだけで、それはきちんと元に戻る。

雨の午後、棚の奥から古い文庫本を一冊取り出した。背は日に焼け、ページの端は少し波打っている。それでも開けば、活字は昔と同じ場所で待っていた。本をスキャンしてデータにすると、この順序が崩れることがある。見開きの左右が入れ替わったり、白紙を一枚飛ばしたせいで、左右の組み合わせが一ページずれたりする。

読みかけのページを覚えておくことも、道具の大切な役目だ。栞を挟むように、閉じたところから開けること。それだけで、本に戻ってくる気持ちの敷居はずいぶん低くなる。文字の大きさや余白も、読みやすさを左右する。ページ全体を収めて眺めたいときもあれば、幅いっぱいに広げて一行ずつ追いたいときもある。

本をスキャンしてデータにすると、この順序が崩れることがある。見開きの左右が入れ替わったり、白紙を一枚飛ばしたせいで、左右の組み合わせが一ページずれたりする。紙の本を読むとき、私たちは右のページから左のページへ、上から下へと目を運ぶ。縦書きの本では、ページをめくる手も右から左へ動く。その身体の動きごと、読書という体験はできあがっている。

読みかけのページを覚えておくことも、道具の大切な役目だ。栞を挟むように、閉じたところから開けること。それだけで、本に戻ってくる気持ちの敷居はずいぶん低くなる。本をスキャンしてデータにすると、この順序が崩れることがある。見開きの左右が入れ替わったり、白紙を一枚飛ばしたせいで、左右の組み合わせが一ページずれたりする。

紙の本を読むとき、私たちは右のページから左のページへ、上から下へと目を運ぶ。縦書きの本では、ページをめくる手も右から左へ動く。その身体の動きごと、読書という体験はできあがっている。ずれた見開きは、読み進めるほどに気持ちが悪い。挿絵の右半分と左半分が別々のページに分かれてしまうこともある。一枚の空白ページを差し込むだけで、それはきちんと元に戻る。

夜の電車で数ページ、昼休みにまた数ページ。細切れの時間でも、前の続きからすぐに読み始められるなら、一冊を読み終えるのは思ったほど難しくない。紙の本を読むとき、私たちは右のページから左のページへ、上から下へと目を運ぶ。縦書きの本では、ページをめくる手も右から左へ動く。その身体の動きごと、読書という体験はできあがっている。

\chapter{ずれた見開き}
紙の本を読むとき、私たちは右のページから左のページへ、上から下へと目を運ぶ。縦書きの本では、ページをめくる手も右から左へ動く。その身体の動きごと、読書という体験はできあがっている。雨の午後、棚の奥から古い文庫本を一冊取り出した。背は日に焼け、ページの端は少し波打っている。それでも開けば、活字は昔と同じ場所で待っていた。

ずれた見開きは、読み進めるほどに気持ちが悪い。挿絵の右半分と左半分が別々のページに分かれてしまうこともある。一枚の空白ページを差し込むだけで、それはきちんと元に戻る。考えてみれば、本の形は長い時間をかけて磨かれてきた道具である。画面の上で読むときも、その形をできるだけ素直に写したい。

文字の大きさや余白も、読みやすさを左右する。ページ全体を収めて眺めたいときもあれば、幅いっぱいに広げて一行ずつ追いたいときもある。夜の電車で数ページ、昼休みにまた数ページ。細切れの時間でも、前の続きからすぐに読み始められるなら、一冊を読み終えるのは思ったほど難しくない。

考えてみれば、本の形は長い時間をかけて磨かれてきた道具である。画面の上で読むときも、その形をできるだけ素直に写したい。考えてみれば、本の形は長い時間をかけて磨かれてきた道具である。画面の上で読むときも、その形をできるだけ素直に写したい。

夜の電車で数ページ、昼休みにまた数ページ。細切れの時間でも、前の続きからすぐに読み始められるなら、一冊を読み終えるのは思ったほど難しくない。読みかけのページを覚えておくことも、道具の大切な役目だ。栞を挟むように、閉じたところから開けること。それだけで、本に戻ってくる気持ちの敷居はずいぶん低くなる。

ずれた見開きは、読み進めるほどに気持ちが悪い。挿絵の右半分と左半分が別々のページに分かれてしまうこともある。一枚の空白ページを差し込むだけで、それはきちんと元に戻る。本をスキャンしてデータにすると、この順序が崩れることがある。見開きの左右が入れ替わったり、白紙を一枚飛ばしたせいで、左右の組み合わせが一ページずれたりする。

ずれた見開きは、読み進めるほどに気持ちが悪い。挿絵の右半分と左半分が別々のページに分かれてしまうこともある。一枚の空白ページを差し込むだけで、それはきちんと元に戻る。紙の本を読むとき、私たちは右のページから左のページへ、上から下へと目を運ぶ。縦書きの本では、ページをめくる手も右から左へ動く。その身体の動きごと、読書という体験はできあがっている。

読みかけのページを覚えておくことも、道具の大切な役目だ。栞を挟むように、閉じたところから開けること。それだけで、本に戻ってくる気持ちの敷居はずいぶん低くなる。考えてみれば、本の形は長い時間をかけて磨かれてきた道具である。画面の上で読むときも、その形をできるだけ素直に写したい。

夜の電車で数ページ、昼休みにまた数ページ。細切れの時間でも、前の続きからすぐに読み始められるなら、一冊を読み終えるのは思ったほど難しくない。考えてみれば、本の形は長い時間をかけて磨かれてきた道具である。画面の上で読むときも、その形をできるだけ素直に写したい。

\chapter{栞のありか}
読みかけのページを覚えておくことも、道具の大切な役目だ。栞を挟むように、閉じたところから開けること。それだけで、本に戻ってくる気持ちの敷居はずいぶん低くなる。紙の本を読むとき、私たちは右のページから左のページへ、上から下へと目を運ぶ。縦書きの本では、ページをめくる手も右から左へ動く。その身体の動きごと、読書という体験はできあがっている。

紙の本を読むとき、私たちは右のページから左のページへ、上から下へと目を運ぶ。縦書きの本では、ページをめくる手も右から左へ動く。その身体の動きごと、読書という体験はできあがっている。文字の大きさや余白も、読みやすさを左右する。ページ全体を収めて眺めたいときもあれば、幅いっぱいに広げて一行ずつ追いたいときもある。

本をスキャンしてデータにすると、この順序が崩れることがある。見開きの左右が入れ替わったり、白紙を一枚飛ばしたせいで、左右の組み合わせが一ページずれたりする。夜の電車で数ページ、昼休みにまた数ページ。細切れの時間でも、前の続きからすぐに読み始められるなら、一冊を読み終えるのは思ったほど難しくない。

本をスキャンしてデータにすると、この順序が崩れることがある。見開きの左右が入れ替わったり、白紙を一枚飛ばしたせいで、左右の組み合わせが一ページずれたりする。考えてみれば、本の形は長い時間をかけて磨かれてきた道具である。画面の上で読むときも、その形をできるだけ素直に写したい。

文字の大きさや余白も、読みやすさを左右する。ページ全体を収めて眺めたいときもあれば、幅いっぱいに広げて一行ずつ追いたいときもある。雨の午後、棚の奥から古い文庫本を一冊取り出した。背は日に焼け、ページの端は少し波打っている。それでも開けば、活字は昔と同じ場所で待っていた。

紙の本を読むとき、私たちは右のページから左のページへ、上から下へと目を運ぶ。縦書きの本では、ページをめくる手も右から左へ動く。その身体の動きごと、読書という体験はできあがっている。夜の電車で数ページ、昼休みにまた数ページ。細切れの時間でも、前の続きからすぐに読み始められるなら、一冊を読み終えるのは思ったほど難しくない。

夜の電車で数ページ、昼休みにまた数ページ。細切れの時間でも、前の続きからすぐに読み始められるなら、一冊を読み終えるのは思ったほど難しくない。夜の電車で数ページ、昼休みにまた数ページ。細切れの時間でも、前の続きからすぐに読み始められるなら、一冊を読み終えるのは思ったほど難しくない。

考えてみれば、本の形は長い時間をかけて磨かれてきた道具である。画面の上で読むときも、その形をできるだけ素直に写したい。考えてみれば、本の形は長い時間をかけて磨かれてきた道具である。画面の上で読むときも、その形をできるだけ素直に写したい。

紙の本を読むとき、私たちは右のページから左のページへ、上から下へと目を運ぶ。縦書きの本では、ページをめくる手も右から左へ動く。その身体の動きごと、読書という体験はできあがっている。紙の本を読むとき、私たちは右のページから左のページへ、上から下へと目を運ぶ。縦書きの本では、ページをめくる手も右から左へ動く。その身体の動きごと、読書という体験はできあがっている。

\chapter{余白と文字}
読みかけのページを覚えておくことも、道具の大切な役目だ。栞を挟むように、閉じたところから開けること。それだけで、本に戻ってくる気持ちの敷居はずいぶん低くなる。考えてみれば、本の形は長い時間をかけて磨かれてきた道具である。画面の上で読むときも、その形をできるだけ素直に写したい。

紙の本を読むとき、私たちは右のページから左のページへ、上から下へと目を運ぶ。縦書きの本では、ページをめくる手も右から左へ動く。その身体の動きごと、読書という体験はできあがっている。雨の午後、棚の奥から古い文庫本を一冊取り出した。背は日に焼け、ページの端は少し波打っている。それでも開けば、活字は昔と同じ場所で待っていた。

読みかけのページを覚えておくことも、道具の大切な役目だ。栞を挟むように、閉じたところから開けること。それだけで、本に戻ってくる気持ちの敷居はずいぶん低くなる。考えてみれば、本の形は長い時間をかけて磨かれてきた道具である。画面の上で読むときも、その形をできるだけ素直に写したい。

読みかけのページを覚えておくことも、道具の大切な役目だ。栞を挟むように、閉じたところから開けること。それだけで、本に戻ってくる気持ちの敷居はずいぶん低くなる。文字の大きさや余白も、読みやすさを左右する。ページ全体を収めて眺めたいときもあれば、幅いっぱいに広げて一行ずつ追いたいときもある。

夜の電車で数ページ、昼休みにまた数ページ。細切れの時間でも、前の続きからすぐに読み始められるなら、一冊を読み終えるのは思ったほど難しくない。雨の午後、棚の奥から古い文庫本を一冊取り出した。背は日に焼け、ページの端は少し波打っている。それでも開けば、活字は昔と同じ場所で待っていた。

考えてみれば、本の形は長い時間をかけて磨かれてきた道具である。画面の上で読むときも、その形をできるだけ素直に写したい。夜の電車で数ページ、昼休みにまた数ページ。細切れの時間でも、前の続きからすぐに読み始められるなら、一冊を読み終えるのは思ったほど難しくない。

本をスキャンしてデータにすると、この順序が崩れることがある。見開きの左右が入れ替わったり、白紙を一枚飛ばしたせいで、左右の組み合わせが一ページずれたりする。紙の本を読むとき、私たちは右のページから左のページへ、上から下へと目を運ぶ。縦書きの本では、ページをめくる手も右から左へ動く。その身体の動きごと、読書という体験はできあがっている。

考えてみれば、本の形は長い時間をかけて磨かれてきた道具である。画面の上で読むときも、その形をできるだけ素直に写したい。雨の午後、棚の奥から古い文庫本を一冊取り出した。背は日に焼け、ページの端は少し波打っている。それでも開けば、活字は昔と同じ場所で待っていた。

ずれた見開きは、読み進めるほどに気持ちが悪い。挿絵の右半分と左半分が別々のページに分かれてしまうこともある。一枚の空白ページを差し込むだけで、それはきちんと元に戻る。読みかけのページを覚えておくことも、道具の大切な役目だ。栞を挟むように、閉じたところから開けること。それだけで、本に戻ってくる気持ちの敷居はずいぶん低くなる。

\end{document}
